\section*{ID9636: Definition of M\_temporal}
\paragraph{Metadata.}
PiID: 4651438407007 | RTEID: RTE:P37418.7279:T5.1059 | UUID: b7f656fa-5bf4-570a-b395-314a28ffdf24

\paragraph{Canonical Statement.}
loor 2\textbackslash{}pi Rn\textbackslash{}Omega/(c\textbackslash{}lambda) \textbackslash{}rfloor\$ distinguishable modes, where \$R\$ is the helix radius, \$n\$ the refractive index, \$\textbackslash{}Omega\$ the rotation frequency, \$c\$ the speed of light, and \$\textbackslash{}lambda\$ the wavelength. For typical parameters \$R = 5\$ millimeters, \$n = 1.45\$, \$\textbackslash{}Omega = 10\textasciicircum{}5\$ revolutions per minute, and \$\textbackslash{}lambda = 500\$ nanometers, this yields approximately \$M\_\{\textbackslash{}text\{OAM\}\} \textbackslash{}approx 10\textasciicircum{}3\$ channels, contributing \$\textbackslash{}log\_2(10\textasciicircum{}3) \textbackslash{}approx 10\$ bits. Polarization contributes an additional \$\textbackslash{}log\_2(2) = 1\$ bit through the two orthogonal polarization states. Temporal multiplexing through rotation provides \$M\_\{\textbackslash{}text\{temporal\}\} = \textbackslash{}lfloor \textbackslash{}Omega T/(2\textbackslash{}pi) \textbackslash{}rfloor\$ time slots within an observation window of duration \$T\$. For \$\textbackslash{}Omega = 10\textasciicircum{}5\$ revolutions per minute and \$T = 1\$ second, this contributes approximately \$M\_\{\textbackslash{}text\{temporal\}\} \textbackslash{}approx 1667\$ slots or \$\textbackslash{}log\_2(1667) \textbackslash{}approx 10.7\$ bits. The total information capacity thus reaches: \textbackslash{}begin\{equation\} C\_\{\textbackslash{}text\{total\}\} = \textbackslash{}log\_2(153600 \textbackslash{}times M\_\{\textbackslash{}text\{OAM\}\} \textbackslash{}times 2 \textbackslash{}times M\_\{\textbackslash{}text\{temporal\}\}) \textbackslash{}approx 39 \textbackslash{}text\{ bits per photon\} \textbackslash{}end\{equation\} This capacity can be distributed across a volumetric display regio

\paragraph{Canonical Equation.}
\[
M_{\text{temporal}} = \lfloor \Omega T/(2\pi) \rfloor
\]

\paragraph{Dictionary.}
Definition of M\_temporal — M\_temporal = lfloor Ω T/(2π) rfloor

\paragraph{Encyclopedia.}
Definition of M\_temporal is an algebraic definition, written M\_temporal = lfloor Ω T/(2π) rfloor. It is an algebraic definition expressing the quantity directly in terms of the variables on its right-hand side. It is built from Ω, π, T, defining M\_temporal. Within the Trawin Zero Point registry it serves as one element of the framework's mathematical narrative.
